\section*{Overview}

Rolling hash turns substrings into fingerprints that can be compared in constant time after preprocessing. It is not a
proof-level string algorithm like KMP or suffix automaton, but it is extremely practical in contests when fast
substring equality is the real need.

\section*{When to Use It}

Use rolling hash when:

\begin{itemize}
  \item many substring-equality checks are needed,
  \item deterministic linear matching is not the main bottleneck,
  \item a tiny collision probability is acceptable or can be reduced with double hashing.
\end{itemize}

\section*{Core Idea}

Interpret the string as a polynomial:
\[
h(s_0 s_1 \dots s_{n-1}) = \sum s_i \cdot base^i \pmod{mod}.
\]
Precompute prefix hashes and powers of the base. Then any substring hash can be extracted in \(O(1)\).

\section*{Key Insight}

The ``rolling'' part is that neighboring substrings differ by one removed character and one added character. Prefix
preprocessing packages that relation so individual substring queries become constant-time arithmetic.

\section*{Operations / Main Technique}

\begin{itemize}
  \item precompute powers and prefix hashes,
  \item query a substring hash in \(O(1)\),
  \item compare two substrings by comparing their normalized hashes.
\end{itemize}

\paragraph{Worked example.}
If I need to compare \(s[l_1 \dots r_1]\) and \(s[l_2 \dots r_2]\) many times, rolling hash turns each comparison into
two prefix-difference computations instead of a character-by-character scan.

\section*{Correctness Intuition}

The hash is designed so concatenation and prefix removal behave algebraically. Taking a prefix difference isolates the
contribution of one substring, and multiplying by the right power aligns hashes from different positions.

The algorithm is probabilistic: equal hashes strongly suggest equal substrings, but collisions are possible.

\section*{Complexity Analysis}

\begin{itemize}
  \item preprocessing: \(O(n)\),
  \item substring hash query: \(O(1)\),
  \item memory: \(O(n)\).
\end{itemize}

\section*{Implementation}

The reference code uses one modulus and one base for clarity. In harder problems I often upgrade that to double hashing
or 64-bit hashing to reduce collision risk.

\section*{Common Pitfalls}

\begin{itemize}
  \item Forgetting that hashing is probabilistic.
  \item Using a poor base or modulus choice.
  \item Comparing unnormalized substring hashes from different positions.
  \item Treating hash equality as a formal proof in adversarial settings.
\end{itemize}

\section*{Variants / Extensions}

\begin{itemize}
  \item Double hashing with two moduli.
  \item 64-bit hashing with natural overflow.
  \item Binary search plus hashing for longest common prefix queries.
\end{itemize}

\section*{Practice Problems}

\begin{itemize}
  \item Many substring-equality queries on one string.
  \item Detect repeated substrings quickly.
  \item Binary-search the longest duplicate or longest common substring length.
\end{itemize}

\section*{References}

\begin{itemize}
  \item \href{https://cp-algorithms.com/string/string-hashing.html}{cp-algorithms: String Hashing}.
  \item \href{https://usaco.guide/gold/hashing?lang=cpp}{USACO Guide: Hashing}.
  \item \href{https://codeforces.com/catalog}{Codeforces Catalog}.
\end{itemize}
